% !TeX root = ../../Book1.tex
% 纲要标题：### 1.4 生成与泛性质
% 纲要位置：计划纲要.md，第 151 行。

\section{生成与泛性质}
\label{b1:sec:0104}

% 纲要内容（仅作写作计划，尚非教材正文）：
%
% * 生成子对象与生成元
% * 自由对象的映射延拓性质
% * 存在性、唯一性与典范同构的区别
%
% ---
%

\subsection{生成子对象与生成元}
\label{b1:subsec:0104:01}

要描述一个很大的代数对象，直接列出所有元素通常并不可行。更有效的办法是先挑选少量元素，再说明从这些元素可以进行哪些运算。生成问题因而分成两个层次：哪些元素能够得到，以及同一个元素有哪些不同的表达式。

\begin{definition}\label{b1:def:generated-monoid}
幺半群 \(M\) 的子集 \(N\) 若包含 \(1_M\) 且对乘法封闭，称为\term[zi-yaobanqun]{子幺半群}[submonoid]。给定 \(S\subseteq M\)，包含 \(S\) 的最小子幺半群称为 \(S\) \term[shengcheng]{生成}[generate]的子幺半群，记为 \(\langle S\rangle_{\mathrm{mon}}\)。若它等于 \(M\)，则称 \(S\) 为一个生成元集合。
\end{definition}
\symidx{\(\langle S\rangle_{\mathrm{mon}}\)：生成子幺半群}

\begin{proposition}\label{b1:prop:generated-monoid}
\(\langle S\rangle_{\mathrm{mon}}\) 存在，既等于所有包含 \(S\) 的子幺半群的交，又等于 \(S\) 中元素的所有有限乘积组成的集合，其中包括空乘积 \(1_M\)。
\end{proposition}
\begin{proof}
包含 \(S\) 的子幺半群至少有 \(M\) 自身。它们的交含有单位元；若两个元素属于交，它们属于每个子幺半群，因而乘积仍属于每一个，也属于交。这给出最小子幺半群。另一方面，所有有限乘积的集合包含 \(S\) 和单位元，并且两个有限乘积的乘积仍为有限乘积；任意包含 \(S\) 的子幺半群都必须包含这些乘积。故两种描述相同。
\end{proof}

例如在 \((\mathbb N,+)\) 中，\(\{2,3\}\) 生成
\(\{0,2,3,4,5,\ldots\}\)。因为 \(1\) 不能写成非负整数系数的 \(2,3\) 之和，它不在其中；每个偶数至少为 \(2\) 时是若干个 \(2\) 的和，每个奇数至少为 \(3\) 时是 \(3\) 加一个偶数。这里 \(6=2+2+2=3+3\) 显示生成表达式一般不唯一。

\begin{proposition}\label{b1:prop:maps-on-generators}
若 \(S\) 生成幺半群 \(M\)，则两个同态 \(f,g:M\rightarrow N\) 在 \(S\) 上相等便在 \(M\) 上相等。
\end{proposition}
\begin{proof}
对非空乘积 \(s_1\cdots s_r\)，同态的值分别为
\(f(s_1)\cdots f(s_r)\) 与 \(g(s_1)\cdots g(s_r)\)，逐因子相等；空乘积的像都是 \(1_N\)。这些元素已穷尽 \(M\)。
\end{proof}

这个命题只保证唯一性，不保证任意指定生成元的像都能形成同态。例如上例中，若想把 \(2\) 送到 \(1\in\mathbb N\)、把 \(3\) 送到 \(1\in\mathbb N\)，则 \(6\) 按两种表达式应同时送到 \(3\) 与 \(2\)，产生矛盾。生成关系正是延拓必须遵守的限制。

\subsection{自由对象的映射延拓性质}
\label{b1:subsec:0104:02}

\begin{definition}\label{b1:def:universal-property}
用对所有目标对象的映射存在性和唯一性来刻画一个构造的条件，称为该构造的\term[fan-xingzhi]{泛性质}[universal property]。例如 \(X^*\) 连同字母映射 \(\iota:X\rightarrow X^*\) 的泛性质是：到任意幺半群的字母映射都能唯一延拓为同态。
\end{definition}

\cref{b1:thm:free-monoid}已经证明这一性质。它同时处理所有目标幺半群，比列举若干可代入的数或矩阵更强。自由对象包含的资料是“对象加指定生成映射”，不是仅仅一个没有标记的幺半群。

\begin{proposition}[生成元与关系]\label{b1:prop:monoid-generators-relations}
设 \(R\) 是 \(X^*\) 中若干对字组成的集合。存在包含这些字对的最小同余关系 \(\equiv_R\)。商幺半群 \(X^*/{\equiv_R}\) 具有如下性质：它到 \(M\) 的同态，恰好对应那些使每对 \((v,w)\in R\) 都满足
\(\widehat u(v)=\widehat u(w)\) 的映射 \(u:X\rightarrow M\)。
\end{proposition}
\begin{proof}
把所有包含 \(R\) 的同余关系相交。这样的关系至少有把所有字视为等价的关系；交仍是等价关系，并且仍与左右拼接相容，故给出最小同余。若 \(u\) 满足所有指定等式，则 \(v\sim w\Longleftrightarrow\widehat u(v)=\widehat u(w)\) 是包含 \(R\) 的同余，因而包含 \(\equiv_R\)。由商幺半群的映射性质，\(\widehat u\) 唯一下降。反过来，商上的同态与商映射复合，在每对指定关系上必取相同值；限制到字母上恢复 \(u\)。这两个过程互为逆过程。
\end{proof}

例如在 \(\{a\}^*\) 中施加 \(a^2=a\)，任何正长度的字都等价于 \(a\)，所以商最多有 \(\varepsilon,a\) 两个元素。它们不会进一步相等：把 \(a\) 送到两元素幺半群 \(\{1,e\}\) 中满足 \(e^2=e\ne1\) 的 \(e\)，可区分二者。于是商正好是这一个两元素幺半群。向任意 \(M\) 指定它的同态，相当于在 \(M\) 中挑选一个满足 \(m^2=m\) 的元素。

\subsection{存在性、唯一性与典范同构的区别}
\label{b1:subsec:0104:03}

构造问题中，“存在”“唯一”和“典范”常被放在一句话里，但它们承担不同职责。存在性要求实际给出对象或论证其可取得；唯一性必须说明相对于哪些固定资料；典范性则强调映射由这些资料强制决定，不依赖额外选择。

\begin{theorem}[泛性质的唯一性]\label{b1:thm:universal-uniqueness}
若幺半群 \(F,F'\) 连同映射 \(i:X\rightarrow F\)、\(i':X\rightarrow F'\) 都满足自由幺半群的映射延拓性质，则存在唯一同构 \(\alpha:F\rightarrow F'\) 满足 \(\alpha i=i'\)。
\end{theorem}
\begin{proof}
由 \(F\) 的泛性质，\(i'\) 唯一延拓为同态 \(\alpha:F\rightarrow F'\)；由 \(F'\) 的泛性质，\(i\) 唯一延拓为同态 \(\beta:F'\rightarrow F\)。有
\((\beta\alpha)i=i\)，而恒等同态也在 \(i\) 上取值为 \(i\)。唯一性因此给出
\(\beta\alpha=\operatorname{id}_F\)。交换两个对象得到
\(\alpha\beta=\operatorname{id}_{F'}\)，所以 \(\alpha\) 是同构。任何满足指定等式的同构首先是延拓 \(i'\) 的同态，故只能是 \(\alpha\)。
\end{proof}

这个唯一同构称为由给定资料确定的\term[dianfan-tonggou]{典范同构}[canonical isomorphism]。它不意味着 \(F\) 与 \(F'\) 是同一个集合，也不意味着任意同构都唯一。例如两个字母的自由幺半群有交换两字母的自同构；只有要求分别保持指定字母时，允许的自同构才只剩恒等映射。

\ProseParagraphBreak
同样的双向构造方法适用于商幺半群和交换消去幺半群的群化：先用一方的泛性质构造映射，再用另一方反向构造，最后由唯一性证明复合为恒等。与猜测两个对象的元素对应相比，这种方法往往更短，也更能说明同构为何自然。使用时仍须先证明有关泛性质；把“显然具有泛性质”当作构造的代替，并不能解决良定义或存在性问题。
